\chapter*{\begin{center}
\vspace{-2.5em}	ВВЕДЕНИЕ \vspace{-1.5em}
\end{center}
}
\addcontentsline{toc}{chapter}{ВВЕДЕНИЕ}

Современные многопользовательские игровые приложения предъявляют высокие требования к точности и надёжности передачи данных между участниками сетевого взаимодействия. Одной из ключевых проблем в данной области является рассинхронизация игровых состояний, то есть несоответствие информации о состоянии игрового мира между клиентами и сервером. Подобные ошибки приводят к снижению качества игрового процесса и ухудшению общего пользовательского опыта.

Целью данной научно-исследовательской работы является сравнительный анализ методов предотвращения рассинхронизации в многопользовательских игровых приложениях.

Для достижения поставленной цели необходимо решить следующие задачи:
\begin{enumerate}
\item провести анализ предметной области синхронизации игровых состояний и передачи данных в многопользовательских играх;
\item формализовать задачу обеспечения согласованности игровых данных между клиентами и сервером;
\item рассмотреть существующие подходы к синхронизации состояния;
\item рассмотреть методы передачи игровых данных;
\item описать типовые игровые данные, подлежащие хранению, обновлению и передаче;
\item определить основные причины возникновения рассинхронизации;
\item сформулировать критерии для оценки и сравнения методов синхронизации;
\item провести сравнительный анализ методов предотвращения рассинхронизации на основе сформулированных критериев.
\end{enumerate}
